\section*{Capítulo \ref{ch:b1:populations} --- Poblaciones y demografía}

\begin{solution}{exo:b1:populations:1}
\omterm{def:b1:populations:population}{Población}: los individuos de una especie en un lugar y un momento, capaces
de cruzarse. Densidad: individuos por unidad de superficie o de volumen.
Dispersión: su patrón espacial --- agregada (una manada, un rodal de
ortigas), uniforme (alcatraces que anidan a un picotazo de distancia,
gobernadoras), aleatoria (dientes de león en un césped).
\end{solution}

\begin{solution}{exo:b1:populations:2}
$l_x$: fracción de una cohorte viva a la edad $x$. $m_x$: hijas por hembra
de edad $x$. $R_0 = \sum l_x m_x$: hijas por hembra recién nacida a lo largo
de su vida. $T = \sum x l_x m_x / R_0$: edad media de las madres.
$R_0 = 1$: cada hembra se reemplaza a sí misma; la \omterm{def:b1:populations:population}{población} es estable.
\end{solution}

\begin{solution}{exo:b1:populations:3}
Tipo III: unos 15 de cada 1000, el \qty{1.5}{\%}. Tipo I: unos 700 de cada
1000, el \qty{70}{\%}.
\end{solution}

\begin{solution}{exo:b1:populations:4}
$\dd N/\dd t = rN$ y $\dd N/\dd t = rN(1 - N/K)$. $r$ es la tasa de
crecimiento per cápita cuando los recursos son ilimitados (nacimientos menos
muertes por individuo y unidad de tiempo); $K$, la \omterm{thm:b1:populations:logistic}{capacidad de carga}, el
tamaño al que el crecimiento se detiene.
\end{solution}

\begin{solution}{exo:b1:populations:5}
$r = \ln(1000)/5 = \qty{1.38}{h^{-1}}$; duplicación, $\ln 2/1.38 =
\qty{0.5}{h}$.
\end{solution}

\begin{solution}{exo:b1:populations:6}
$l_x m_x = 0, 0.5, 0.6, 0.2, 0$: $R_0 = 1.3$. $T = (1\times 0.5 +
2\times 0.6 + 3\times 0.2)/1.3 = 2.3/1.3 = 1.77$. $r = \ln 1.3/1.77 =
\qty{0.15}{}$ por unidad de tiempo: crece.
\end{solution}

\begin{solution}{exo:b1:populations:7}
$N = 50\times 70/7 = 500$. Solo quedaban 40 peces marcados: $N = 40\times
70/7 = 400$; la primera estimación era demasiado alta en una cuarta parte
(la fracción marcada de la \omterm{def:b1:populations:population}{población} era menor de lo supuesto).
\end{solution}

\begin{solution}{exo:b1:populations:8}
$0.2\times 100\times 0.9 = 18$; $0.2\times 500\times 0.5 = 50$;
$0.2\times 900\times 0.1 = 18$ al año. Máximo $rK/4 = 50$ al año, en
$N = 500$.
\end{solution}

\begin{solution}{exo:b1:populations:9}
Diente de león: $r$ (miles de semillas, anual). Elefante: $K$ (una cría cada
cuatro años, vida larga, cuidados). Bacalao: $r$ (millones de huevos, sin
cuidados). Gorila: $K$ (una cría cada cuatro años, décadas de vida). Mosca
doméstica: $r$ (cientos de huevos, madura en diez días).
\end{solution}

\begin{solution}{exo:b1:populations:10}
$r = \ln(29/10)/2 = \qty{0.53}{h^{-1}}$; $K \approx 660$. A las \qty{8}{h}:
$660/(1 + 65e^{-4.24}) = 660/(1 + 0.94) = 340$, frente a 351 contadas.
Máximo en $K/2 = 330$, alcanzado en $t = \ln 65/0.53 = \qty{7.9}{h}$,
donde $\dd N/\dd t = rK/4 = 87\times 10^5$ \omterm{def:b1:cell-unit-of-life:cell}{células} por mililitro y
hora.
\end{solution}

\begin{solution}{exo:b1:populations:11}
La regulación exige que la tasa per cápita caiga al subir $N$, de modo que
$N$ vuelva hacia un equilibrio tras una perturbación; un factor que mata una
fracción fija cambia $N$ pero no la dependencia de la tasa respecto de $N$,
así que después la \omterm{def:b1:populations:population}{población} reanuda el mismo crecimiento. El tiempo limita
a los insectos golpeando con suficiente frecuencia, y al azar, como para que
la \omterm{def:b1:populations:population}{población} quede rezagada antes de poder acercarse a techo alguno: una
\omterm{def:b1:populations:population}{población} no regulada, mantenida baja por perturbaciones repetidas, cuyo
tamaño medio depende de la frecuencia de las malas temporadas y no de una
\omterm{thm:b1:populations:logistic}{capacidad de carga}.
\end{solution}

\begin{solution}{exo:b1:populations:12}
Dos poblaciones del mismo tamaño, una con una base joven ancha y otra con
clases de edad parejas, tienen futuros distintos: la primera crecerá durante
décadas aunque su $m_x$ caiga ahora mismo al nivel de reemplazo, porque sus
muchos jóvenes todavía no se han reproducido (inercia demográfica); la
segunda, no. $N$ registra el pasado; la distribución de edades y el
calendario de $m_x$ sobre ella son aquello sobre lo que actúan las
ecuaciones, y la pirámide predice el tamaño de la generación siguiente,
mientras que $N$ no dice nada al respecto.
\end{solution}

\begin{solution}{pb:b1:populations:1}
\textbf{1.} $r = \ln(29/10)/2 = \qty{0.53}{h^{-1}}$.
\textbf{2.} $\ln 2/0.53 = \qty{1.3}{h}$.
\textbf{3.} $K \approx 660\times 10^5 = \qty{6.6e7}{}$ \omterm{def:b1:cell-unit-of-life:cell}{células} por
mililitro.
\textbf{4.} $N(t) = 660/(1 + 65e^{-0.53t})$: a las 4 h, $660/(1 + 7.8) =
75$ (71 contadas); a las 8 h, 340 (351); a las 12 h, $660/(1 + 0.112) = 594$
(595).
\textbf{5.} En $K/2 = 330$, alcanzado en $t = \ln 65/0.53 = \qty{7.9}{h}$;
tasa máxima $rK/4 = 0.53\times 660/4 = 87\times 10^5$ \omterm{def:b1:cell-unit-of-life:cell}{células} por
mililitro y hora.
\textbf{6.} $10e^{0.53\times 12} = 10\times 578 = 5780$: diez veces las 595
observadas.
\textbf{7.} El agotamiento del azúcar (o del oxígeno); la acumulación de
etanol y de ácidos que envenenan las \omterm{def:b1:cell-unit-of-life:cell}{células}.
\textbf{8.} $q_0 = 0.40$; $q_4 = 1 - 0.34/0.40 = 0.15$: mortalidad temprana
intensa y después un ritmo más constante --- entre el tipo II y el III, como
en la mayoría de los grandes mamíferos silvestres.
\textbf{9.} $l_x m_x = 0, 0, 0.312, 0.368, 0.320, 0.272, 0.182, 0.080,
0.018, 0$: $R_0 = 1.55$.
\textbf{10.} $\sum x l_x m_x = 0.624 + 1.104 + 1.280 + 1.360 + 1.092 +
0.560 + 0.144 = 6.16$; $T = 6.16/1.55 = \qty{4.0}{yr}$.
\textbf{11.} $r = \ln 1.55/4.0 = \qty{0.11}{yr^{-1}}$; duplicación en
\qty{6.3}{yr}.
\textbf{12.} $200e^{1.1} = 600$ ciervos.
\textbf{13.} $N(10) = 600/(1 + 2e^{-1.1}) = 600/1.67 = 360$; pasa de 500
cuando $(K - N_0)/N_0\,e^{-rt} = 0.2$, es decir, $e^{-rt} = 0.1$, $t = \ln
10/0.11 = 21$ años.
\textbf{14.} $120\times 150/18 = 1000$ percas.
\textbf{15.} 100 marcadas: $100\times 150/18 = 833$; la primera estimación
era demasiado alta (había menos peces marcados sueltos de lo supuesto).
\textbf{16.} $rK/4 = 100$ al año, en $N = 500$.
\textbf{17.} $T_2 = \ln 2/0.4 = \qty{1.7}{yr}$.
\textbf{18.} Captura $hN$ con $h = 0.2$: $rN(1 - N/K) = hN$ da $N^* =
K(1 - h/r) = 1000\times 0.5 = 500$ y una captura de \num{100} al año ---
por debajo del rendimiento máximo sostenible, pero una fracción se ajusta
sola a la existencia: si esta se reduce a la mitad, también la captura, y la
\omterm{def:b1:populations:population}{población} se recupera.
\textbf{19.} En 500: $0.4\times 500\times 0.5 = 100 < 120$: la \omterm{def:b1:populations:population}{población}
decae. En 300: $0.4\times 300\times 0.7 = 84 < 120$: decae más deprisa ---
una vez por debajo de $K/2$, una cosecha fija superior al rendimiento máximo
la lleva a la extinción.
\textbf{20.} En $K/2$ el excedente es máximo, pero cualquier error ---
una \omterm{def:b1:populations:population}{población} sobreestimada en una quinta parte, un mal año que recorte $r$
--- convierte una extracción sostenible en un declive, y por debajo de $K/2$
el excedente mengua a la vez que la \omterm{def:b1:populations:population}{población}: el equilibrio es inestable
frente a la sobreexplotación. Cosechar algo por encima de $K/2$ deja margen.
\textbf{21.} $R_0$ se reduce a la mitad, $0.78 < 1$: la manada mengua hasta
caer por debajo de 400 y $m_x$ se recupera --- un factor dependiente de la
densidad, que regula la manada cerca de 400.
\textbf{22.} No: retira el \qty{40}{\%} sea cual sea el tamaño y deja
inalterada la tasa per cápita; en los años siguientes la manada vuelve a
crecer con la misma $r$ desde un punto de partida más bajo --- una
perturbación, no una regulación.
\textbf{23.} $0.11N(1 - N/600) = 30$: $N^2 - 600N + 163\,600 = 0$,
$N = 300\pm\sqrt{90\,000 - 163\,600}$ --- sin solución real: 30 al año
supera el excedente máximo $rK/4 = 16.5$, y los lobos llevan la manada a la
extinción. Con 15 al año: $N = 300\pm\sqrt{90\,000 -
81\,800} = 300\pm 91$, es decir, 209 y 391; el superior es estable
(por encima de él el excedente es menor que 15 y la manada retrocede; por
debajo, mayor, y sube), y el inferior, inestable.
\textbf{24.} Los efectivos del lince responden al alimento con retraso: una
subida de las liebres eleva los nacimientos y la supervivencia de los linces
a lo largo del año o los dos siguientes, y un desplome de las liebres solo
mata de hambre a los linces cuando se han acabado su \omterm{def:b1:lipids:triglyceride}{grasa} y los cachorros
de la temporada.
\textbf{25.} Levadura: $K = \qty{6.6e7}{}$ \omterm{def:b1:cell-unit-of-life:cell}{células} por mililitro, $r =
\qty{0.53}{h^{-1}}$. Ciervos: $R_0 = 1.55$, $T = \qty{4.0}{yr}$, $r =
\qty{0.11}{yr^{-1}}$.
\end{solution}
